\documentclass[12pt,a4paper]{article}

\usepackage[T2A]{fontenc}
\usepackage[utf8]{inputenc}
\usepackage[russian]{babel}
\usepackage{geometry}
\usepackage{enumitem}
\usepackage{listings}
\usepackage{xcolor}
\usepackage{hyperref}
\usepackage{longtable}
\usepackage{array}

\geometry{left=25mm,right=15mm,top=20mm,bottom=20mm}
\setlength{\parindent}{1.25cm}
\setlength{\parskip}{0pt}

\lstset{
    basicstyle=\ttfamily\small,
    frame=single,
    breaklines=true,
    showstringspaces=false,
    columns=fullflexible,
    keepspaces=true
}

\begin{document}
\begin{titlepage}
\begin{center}
Федеральное государственное автономное образовательное учреждение\\
высшего образования\\[2mm]
\textbf{«Национальный исследовательский университет\\
``Высшая школа экономики''»}\\[3mm]
Факультет компьютерных наук\\
Департамент программной инженерии
\vfill
{\Large\bfseries Отчёт по индивидуальному заданию № 2}\\[4mm]
{\large по дисциплине «Архитектура компьютера и операционные системы»}\\[7mm]
{\Large\bfseries Вариант 14. «Задача об умных цветах»}\\[3mm]
{\large Многопоточная имитационная программа}
\vfill
\begin{flushright}
Выполнила: Гончарова Екатерина\\
Группа: БПИ-253
\end{flushright}
\vfill
Москва, 2026
\end{center}
\end{titlepage}

\tableofcontents
\newpage

\section{Постановка задачи}

Требуется разработать консольное приложение на языке C для ОС Linux,
имитирующее взаимодействие умных цветов и садовников с помощью
\textbf{независимых потоков исполнения}. По условию варианта 14 каждый
цветок самостоятельно изменяет свою влажность, определяет необходимость
полива и посылает заявку. Садовники реагируют на заявки и конкурируют за
право их исполнения. Для одной заявки может быть назначен только один
садовник. Повторный полив в течение защитного интервала вызывает перелив.

Состояния цветка: нормальное развитие (\texttt{NORMAL}), увядание
(\texttt{WILTING}), полив (\texttt{WATERING}), восстановление
(\texttt{RECOVERING}), засыхание (\texttt{DRY}), гибель от перелива
(\texttt{DEAD\_OVERWATERED}). Приложение обязано отображать смену
состояний, создание и снятие заявки, её принятие садовником, конкуренцию
за заявки, полив, восстановление, засыхание и перелив.

Для реализации используются \texttt{pthread}, мьютексы и семафоры.
Каждый участник модели выполняет собственный цикл обработки событий.

\section{Соответствие сущностей потокам}

Создаются \(N\) потоков цветов и \(M\) потоков садовников, то есть всего
\(N+M\) рабочих потоков. Главный поток используется только для
создания и завершения рабочих потоков, отслеживания времени и вывода
итоговых данных.

\subsection{Поток цветка}

Каждый поток получает указатель на общую модель и индекс только
своего цветка. Через условный такт он:
\begin{enumerate}[noitemsep]
\item захватывает мьютекс модели;
\item изменяет влажность и состояние своего цветка;
\item при достижении порога создаёт одну активную заявку;
\item отправляет уведомление через семафоры садовникам;
\item освобождает мьютекс и ждёт следующий такт.
\end{enumerate}

Поток не создаёт новую заявку до завершения или отмены предыдущей.
Засохший либо погибший цветок больше не меняется и не подаёт заявки.

\subsection{Поток садовника}

Каждый садовник ждёт на собственном семафоре и не загружает процессор
циклом опроса. После получения уведомления он под мьютексом выбирает
свободную заявку по правилу \texttt{first} или \texttt{driest}.
При успешном назначении запоминает индекс цветка и начинает работать.
При отсутствии свободной заявки, если другая заявка уже принята,
фиксируется конфликт. Два потока, конкурирующие за заявку, не могут
успешно назначить её одновременно: проверка доступности и назначение
выполняются в критической секции.

Движение садовника и полив занимают заданное число шагов времени.
Садовник спит между этими шагами, не удерживая мьютекс. Во время
подхода цветок может окончательно засохнуть; садовник тогда прекращает
обслуживание. При завершении полива восстановление влажности и снятие
заявки выполняются атомарно.

\subsection{Главный поток}

Главный поток создаёт необходимые объекты синхронизации и рабочие
потоки с помощью \texttt{pthread\_create()}, после чего ожидает истечения
заданного периода наблюдения либо прекращения активности всех цветов.
Он завершает работу при \texttt{Ctrl+C} и ожидает каждый запущенный поток
через \texttt{pthread\_join()}.

\section{Синхронизация и общая память}

\subsection{Мьютекс состояния}

\texttt{model\_mutex} защищает массив цветов, массив садовников, признаки
активных заявок и назначенных садовников, статистические счётчики и
флаг завершения. Ни один рабочий поток не читает и не изменяет эти
общие переменные вне критической секции. В частности, принятие заявки
осуществляется как неделимая операция.

Упрощённая схема выбора и назначения заявки:
\begin{lstlisting}[language=C]
pthread_mutex_lock(&simulation->model_mutex);
update_gardener(gardener, simulation->flowers,
                simulation->settings, simulation->current_tick);
pthread_mutex_unlock(&simulation->model_mutex);
\end{lstlisting}

Вызов \texttt{update\_gardener} при свободном садовнике находит доступную
заявку, проверяет \texttt{request\_active}, \texttt{assigned\_gardener}
и состояние цветка, после чего закрепляет её за одним садовником.

\subsection{Семафоры и уведомления}

У каждого садовника есть свой семафор уведомлений.
Когда появляется заявка, поток цветка вызывает \texttt{sem\_post()}
для каждого садовника. Свободный садовник ожидает сигнал с помощью
\texttt{sem\_wait()}. После пробуждения он пытается принять заявку. Таким образом два садовника могут реагировать на одну и ту же
заявку, но только один получает право на полив.

Уведомление не эквивалентно гарантированному назначению: заявка к
моменту пробуждения может быть уже принята или отменена. Накопившиеся
уведомления безопасны, так как реальное состояние всегда проверяется
под мьютексом.

\subsection{Потокобезопасный вывод}

\texttt{log\_mutex} сериализует запись событий в стандартный вывод и
\path{simulation.log}. Используются \texttt{open()}, \texttt{write()} и
\texttt{close()}; вспомогательная функция повторяет запись при частичном
\texttt{write()} и при \texttt{EINTR}. Это предотвращает перемешивание
фрагментов разных сообщений.

\section{Входные параметры и пользовательский интерфейс}

Интерфейс поддерживает интерактивный ввод и запуск с файлом параметров. Программа принимает десять
исходных параметров, соответствующих условию: количество цветов,
массив индивидуальных скоростей высыхания, порог заявки, порог
засыхания, длительность полива, восстановления, защитного интервала,
число садовников, скорость их работы и правило выбора заявки.

Возможны два способа запуска: интерактивный, без аргументов, и по
конфигурационному файлу \texttt{key=value}. Например:
\begin{lstlisting}
make
./build/smart_flowers
./build/smart_flowers tests/balanced.cfg
make check
\end{lstlisting}

Модель использует 50 условных тактов с паузой 25 миллисекунд между
шагами каждого потока. Длительность такта и число тактов определены
именованными константами \texttt{TICK\_DURATION\_MS} и
\texttt{OBSERVATION\_PERIOD}. Это параметры визуализации и границы
наблюдения; они не изменяют формат входных конфигураций.

\section{Корректное завершение}

Программа заканчивается после 50 условных тактов, при прекращении
активности всех цветов либо при получении \texttt{SIGINT}. Потоки
наследуют заблокированный \texttt{SIGINT}; поэтому сигнал обрабатывается
главным потоком. Обработчик меняет только \texttt{sig\_atomic\_t}.
Главный поток устанавливает защищённый мьютексом флаг остановки,
пробуждает ожидающих садовников, дожидается завершения всех потоков,
печатает статистику и удаляет примитивы синхронизации.

Незаконченный к моменту остановки полив не включается в число
\emph{успешно завершённых} поливов. В итоговой статистике выводятся
число живых, засохших и погибших от перелива цветов, влажность
каждого цветка и число завершённых поливов. В журнале дополнительно
выводятся количество созданных заявок и конфликтов.

\section{Проверка программы}

\subsection{Тестовые сценарии}

Предусмотрено десять конфигураций с разными сочетаниями числа потоков,
порогов, длительности полива и восстановления, скорости садовников и
правил выбора заявок:
\begin{itemize}[noitemsep]
\item \texttt{balanced.cfg}: обычная работа двух садовников;
\item \texttt{driest\_policy.cfg}: выбор наиболее сухого цветка;
\item \texttt{single\_gardener.cfg}: очередь при одном садовнике;
\item \texttt{overwatering.cfg}: проверка гибели от перелива;
\item \texttt{fast\_drying.cfg}: быстрое засыхание;
\item \texttt{many\_gardeners.cfg}: конкуренция нескольких садовников;
\item \texttt{large\_garden.cfg}: 25 цветов и 10 садовников;
\item \texttt{long\_recovery.cfg}: длительное восстановление;
\item \texttt{no\_protection.cfg}: нулевой защитный интервал;
\item \texttt{one\_flower.cfg}: граничный случай с одним цветком.
\end{itemize}

Команда \texttt{make check} запускает все десять конфигураций с помощью
небольшого скрипта на Python. Скрипт проверяет код завершения программы
и наличие итоговой статистики. Для \texttt{overwatering.cfg} дополнительно
проверяется, что хотя бы один цветок погиб от перелива.

Эти проверки помогают находить ошибки запуска и очевидные нарушения
логики моделирования. Они не проверяют все возможные порядки выполнения
потоков. При ручной проверке можно запустить отдельную конфигурацию и
посмотреть журнал \path{simulation.log}. Порядок событий в нём
может отличаться между запусками.

\subsection{Инварианты}

При работе программы должны соблюдаться следующие условия:
\begin{enumerate}[noitemsep]
\item одну активную заявку получает не более одного садовника;
\item каждый садовник обслуживает не более одного цветка;
\item один цветок имеет не более одной активной заявки;
\item закрытая или отменённая заявка не может быть снова обслужена;
\item погибший цветок не формирует заявки и не продолжает активную работу.
\end{enumerate}

\section{Особенности реализации}

Каждый цветок и садовник работает в отдельном потоке. Для доступа к
общему состоянию используются мьютексы, для уведомлений --- семафоры.
Программа использует один мьютекс для состояния модели. Это упрощает
назначение заявок: проверка и закрепление выполняются вместе.

Общий мьютекс упрощает код, но при большом количестве потоков может
ограничивать скорость работы. Перемещение садовника задаётся количеством
работы, а не координатами. Также у потоков нет общего барьера в конце
каждого такта: они выполняют шаги независимо друг от друга. Поэтому
записи с одинаковым номером такта не обязательно относятся к одному
моменту времени.

\section{Итог}

Разработана многопоточная модель умных цветов и садовников
с конфигурационным вводом, журналом и итоговой статистикой. Каждый цветок и каждый садовник имеет отдельный
поток, а общий доступ синхронизируется мьютексами и семафорами.
Предусмотрено нормальное и досрочное завершение, ведение журнала и
автоматизированная проверка.

\end{document}
